\section{Introduction}
Phase-field simulation is widely used to study the formation and evolution of material microstructures \citep{moelans2008,tourret2022}. A grain or phase is represented by a continuous field that changes smoothly across an interface. The interface moves as the fields evolve, so its position need not be tracked explicitly. This description is useful when boundaries curve, meet or disappear during solidification, grain growth, coarsening and solid-state transformation. It also permits interface motion to be coupled to other processes: solidification models, for example, combine the phase fields with heat and solute transport \citep{boettinger2002}.

In a polycrystal, multiple fields distinguish the individual grains and allow their shapes, sizes and topology to evolve within one calculation \citep{krill2002}. Multiphase-field formulations accommodate different interfacial energies and mobilities for the boundaries between them \citep{steinbach1996,steinbach1999,steinbach2009}. These boundaries are connected through junctions. Their motion must therefore respect both interface kinetics and interfacial force balance, rather than reproduce the speed of an isolated planar boundary alone. Classical curvature-driven grain-boundary motion provides one sharp-interface reference \citep{mullins1956}; bulk driving forces introduce an additional contribution to interface motion. An error near a junction can consequently affect the shape and disappearance time of an entire neighboring grain.

Nanocrystalline materials provide context because their grain structures and stability change during thermal exposure. Studies of triple-junction drag show that finite junction mobility can affect grain-growth rates, size distributions and the transition from junction- to boundary-controlled growth \citep{johnson2014,miyoshi2026}. The direct numerical motivation here is strongly driven transformation-front motion: junction errors at practical diffuse widths can distort adjoining boundary shapes and grain-disappearance times \citep{hoffrogge2025}. Physical nanoscale junction drag is a distinct mechanism; the prescribed sharp-interface tests below neither calibrate that mobility nor validate a particular nanomaterial.

Computational interfaces are often widened to reduce simulation cost, so a quantitative method must reproduce intended motion at its chosen width as well as resolve its equations. Thin-interface analysis and corrections have been particularly successful for solidification \citep{karma1998,karma2001,eiken2006}; narrowing the region of latent-heat release or solute redistribution can further suppress anomalous interface effects there \citep{kim2004}. At multiphase junctions, several profiles and force channels coexist, and correct planar-front motion does not ensure correct junction motion. Benchmarks show sensitivity to interfacial-energy terms, mobility interpolation, resolution and bulk-force allocation \citep{daubner2023,nani2024,hoffrogge2025}. Speed, whole-boundary shape and event timing are therefore complementary measurements.

Several methods already address parts of this problem. Grose and Asimow \citep{grose2022} redistribute bulk driving forces among phase fields and combine this with profile correction and recursive transfer at junctions. Feyen and Moelans \citep{feyen2023} separate normal-profile restoration from physical curvature weighting in a high-driving-force model. Zhang et al. \citep{zhang2023large} extend the driving-force term to remain constant along interface normals in driven Allen--Cahn and concentration-coupled models, a different formulation from the local capillary modification tested here. Profile-informed finite-difference correction \citep{eiken2012}, isotropic spatial operators \citep{tang2024}, and hybrid nonlinear preconditioning for double-obstacle potentials \citep{sakane2026} reduce other important errors. Improved time integration can reduce the work needed to follow a trajectory \citep{seiz2026}, but does not by itself remove an inaccurate finite-width force balance.

We examine a restricted modification: multiply the existing capillary contribution by a factor determined from nearby ternary coexistence, retain the bulk-force law, and extend the signal beyond the ternary core to adjoining binary-interface segments. A local maximum filter implements this construction without interface-normal reconstruction. It modifies the finite-width model toward a specified sharp-interface target; it is not a more accurate integrator of the unchanged equations.

The question is whether this spatial prescription improves driven unequal-energy junction shapes and kinetics at fixed numerical settings. A matched comparison isolates its effect; cold interactions test whole-boundary correction and optional persistent suppression after correction neighborhoods overlap; parameter perturbations assess sensitivity to empirical strength and radius. We also insert the fixed prescription into OpenPhase \citep{tegeler2017,openphase}, retaining its native force assembly and constraint handling. The contribution is this spatial source construction and its tested whole and history-dependent forms. A source-based Grose--Asimow reconstruction provides a scoped comparator, not an exhaustive optimization. These benchmarks assess finite-width numerical accuracy and do not select a universal correction.
